% !PRINTED 559
\weberpage{559}
\begin{center}
\gesperrt{Eighteenth Chapter.}

\gesperrt{Discriminants.}

\rule{0.3\textwidth}{0.4pt}

\S.\,155.

\gesperrt{Minimum of a Quadratic Form.}
\end{center}

On an entirely new path \gesperrt{Minkowski} has proved the finiteness of the number of classes\footnote{Ueber die positiven quadratischen Formen etc. \gesperrt{Crelle}, Bd.\ 107, 1891. More thoroughly still in the book by \gesperrt{Minkowski}, ``Geometrie der Zahlen'', Leipzig 1896.}, a path which starts from peculiar considerations, capable of numerous applications, and in doing so leads to the proof, long sought in vain, of a theorem on the field discriminant. We shall here communicate briefly the principal points of these investigations.

\gesperrt{Minkowski} starts from the consideration of a quadratic form of $n$ variables

\weq{1}{f(x_1, x_2, \ldots, x_n) = \Sigma^{i,k} a_{i,k}\, x_i x_k}

with real coefficients $a_{i,k}$, of which it is assumed that it can be decomposed into $n$ and no fewer positive squares of linear functions (a positive form, Bd.\ I, \S.\,81 to 83). Let one of these representations be

\weq{2}{f(x_1, x_2, \ldots, x_n) = \xi_1^2 + \xi_2^2 + \cdots + \xi_n^2,}
wherein
\weq{3}{\xi_i = \alpha_{1,i} x_1 + \alpha_{2,i} x_2 + \cdots + \alpha_{n,i} x_n}

is a system of linear functions of the $x_i$ with non-vanishing determinant. If we set

\weq{4}{A = \Sigma \pm \alpha_{1,1}\alpha_{2,2}\ldots\alpha_{n,n},}
%CONT>>
